\section{Detailed Related Work}
\label{appendix:more_related_work}

\textbf{Boosting LLM's Reasoning Capacity}. 
Various approaches have been explored to strengthen the reasoning abilities of LLMs. Several methods focus on using high-quality data for training. For instance, RFT \cite{yuan2024advancingllmreasoninggeneralists, yuan2023scalingrelationshiplearningmathematical} fine-tunes the pretrained base model using correct samples generated by a supervised fine-tuned model. Similarly, RestEM \cite{singh2024humandatascalingselftraining} adopts an iterative self-training strategy, repeatedly retraining the base model on high-quality reasoning traces produced by previously obtained checkpoints. Preference-based methods like DPO \cite{rafailov2024directpreferenceoptimizationlanguage,rafailov2024rqlanguagemodel} optimize models by contrasting correct and incorrect reasoning outputs. Search-guided approaches, like Monte Carlo Tree Search (MCTS) \cite{zhang2024restmctsllmselftrainingprocess, chen2024alphamathzeroprocesssupervision}, help models discover improved reasoning paths during inference. Another significant direction involves RL frameworks, which typically adopt policy gradient-based methods and differ mainly in advantage estimation granularity. For example, GRPO \cite{shao2024deepseekmathpushinglimitsmathematical} and RLOO \cite{ahmadian2024basicsrevisitingreinforcestyle} estimate trajectory-level advantages, while PPO \cite{schulman2017proximalpolicyoptimizationalgorithms, ouyang2022traininglanguagemodelsfollow} estimates token-level advantages using a dedicated critic model. Recent work \cite{kazemnejad2024vineppounlockingrlpotential} highlights that accurately training the critic in PPO is challenging, and proposes VinePPO, a critic-free alternative that partitions the reasoning trajectory into discrete steps using heuristic rules (e.g., line breaks) and estimates step-level advantages through Monte Carlo (MC) sampling. PRIME \cite{cui2025processreinforcementimplicitrewards} simultaneously trains a process reward model using outcome rewards during policy training, providing process rewards to better guide policy optimization, while our work does not rely on reward models. Most recently, several works have proposed improvements upon the original GRPO method, like DAPO \cite{yu2025dapoopensourcellmreinforcement} and Dr. GRPO \cite{liu2025understandingr1zeroliketrainingcritical}, based on the trajectory-level advantage estimation.

\textbf{Fine-Grained Reward Signals}. 
A common approach in prior work assigns a single binary reward to the final output of LLMs by comparing it against the ground truth. However, this sparse reward provides limited guidance, resulting in the difficulty of credit assignment during training. To address this challenge, \cite{lightman2023letsverifystepstep} initially proposed ``Process Reward'', which involves manually judging correctness at every intermediate reasoning step. Subsequent works, such as \cite{wang2024mathshepherdverifyreinforcellms}, \cite{luo2024improvemathematicalreasoninglanguage}, and \cite{ma2023letsrewardstepstep}, automated the process reward generation without manual labeling through rollouts. These methods focus on training an external Process Reward Model (PRM) to provide intermediate rewards for policy optimization and enable Best-of-N (BoN) selection during inference. In practice, when applying PRM for policy optimization, existing approaches typically combine step-level rewards with the final binary correctness rewards and still employ standard RL algorithms such as PPO or GRPO for training. In contrast, our approach differs fundamentally by improving the RL optimization algorithm itself, rather than relying on external reward models.  We introduce an MC-based estimation method to directly compute segment-level advantages from the current policy, aiming at potentially reducing gradient estimation variance and enabling finer-grained credit assignment during optimization. This strategy helps stabilize training and leads to more effective policy updates compared to using sparse, trajectory-level rewards alone.  Notably, our MC advantage estimation framework remains fully compatible with the PRM approach: intermediate process reward signals generated via PRMs could be integrated with our trajectory evaluations. Such integration may combine the strengths of both methodologies, providing further potential improvements in overall model performance. In Appendix \ref{appendix:incorporate_process_reward}, we provide an analysis of integrating the process reward into our framework.